\section{Working routine}
\label{sec:routine}

The installation is done once. This section describes what is done every time: how a working
session begins, how results are shared with the team, and how a session ends.

\subsection{Begin a session}

\step{Open the project}
Open a new PowerShell window (Step~\ref{st:openps}) and give the opening sequence.
\begin{ps}[Opening sequence, command 1 of 2]
cd C:\ee105\camera-to-latex
\end{ps}
\begin{ps}[Opening sequence, command 2 of 2]
.\.venv\Scripts\Activate.ps1
\end{ps}
\checkpoint{The prompt reads:}
\venvprompt

\step{Fetch the team's latest changes}\label{st:pull}
\begin{ps}
git pull
\end{ps}
\begin{out}[What you should see when nothing has changed]
Already up to date.
\end{out}
When team members have uploaded changes since your last session, Git instead lists the files it
updated. Begin every session with this command; work done on an outdated copy is the usual
cause of the conflicts described below.
\checkpoint{The output is \emph{Already up to date} or a list of updated files, and it contains
neither \emph{error} nor \emph{CONFLICT}.}

\subsection{Share your work}

Sharing a change takes three commands: \cmd{git add} selects the files, \cmd{git commit} records
them on your laptop with a describing sentence, and \cmd{git push} uploads the record to GitHub.
The companion document gives these commands in full, with the real file names, at every point
where something is to be shared. Their general form is shown here.

\step{See what has changed}
\begin{ps}
git status
\end{ps}
\begin{out}[What you should see when you have changed the file config.py]
On branch main
Your branch is up to date with 'origin/main'.

Changes not staged for commit:
        modified:   config.py
\end{out}
\checkpoint{Every file listed is one that you meant to change. If a file is listed that you did
not mean to change, do not continue; ask the lead.}

\step{Select, record and upload}
In the first command, replace the placeholder by the name of the changed file as \cmd{git
status} printed it. In the second, replace the placeholder by a short sentence that says what
you did; keep the quotation marks.
\begin{center}
\texttt{git add }\ph{FILE}\\[4pt]
\texttt{git commit -m "}\ph{MESSAGE}\texttt{"}\\[4pt]
\texttt{git push}
\end{center}
\begin{out}[Example: sharing a change of config.py]
(.venv) PS C:\ee105\camera-to-latex> git add config.py
(.venv) PS C:\ee105\camera-to-latex> git commit -m "Set camera orientation"
[main 3f2c1ab] Set camera orientation
 1 file changed, 2 insertions(+), 2 deletions(-)
(.venv) PS C:\ee105\camera-to-latex> git push
To https://github.com/ashahvos/camera-to-latex.git
   a1b2c3d..3f2c1ab  main -> main
\end{out}
\checkpoint{The output of \cmd{git push} ends with a line of the form \texttt{main -> main}.
The change is now visible on the repository's page on GitHub.}

Three rules keep the shared repository in order.
\begin{enumerate}
\item Name the files you add. Do not use commands that add ``everything''; they sweep up files
  that were never meant to be shared.
\item Write messages that a team member will understand in a week: \emph{Add photos of test
  lines T01 to T60}, not \emph{update}.
\item Never share passwords, and never try to add the folders \file{.venv}, \file{data} or
  \file{out}; Git ignores them for good reason.
\end{enumerate}

\subsection{When Git objects}

\step{If the upload is rejected}
\begin{out}[The message]
 ! [rejected]        main -> main (fetch first)
error: failed to push some refs to 'https://github.com/ashahvos/camera-to-latex.git'
\end{out}
Someone uploaded a change after your last \cmd{git pull}. Nothing is lost. Fetch their change
and upload again:
\begin{ps}
git pull
\end{ps}
\begin{ps}
git push
\end{ps}
If Notepad opens during \cmd{git pull} with a sentence that begins \emph{Merge branch}, Git is
proposing a description for the combination of the two changes. The proposal is adequate: close
Notepad (\menu{File}\then\menu{Exit}), and the command completes.

\step{If Git reports a conflict}
\begin{out}[The message]
CONFLICT (content): Merge conflict in config.py
Automatic merge failed; fix conflicts and then commit the result.
\end{out}
You and a team member changed the same lines of the same file, and Git cannot decide between
the two versions. Return the project to its state before the pull with
\begin{ps}
git merge --abort
\end{ps}
and send the lead a screenshot. Your own change is still recorded on your laptop; the lead will
say which version should prevail.
\checkpoint{\cmd{git status} no longer mentions a merge.}

\step{If Git refuses to pull because of local changes}
\begin{out}[The message]
error: Your local changes to the following files would be overwritten by merge:
\end{out}
You changed a file without recording the change, and the same file has changed on GitHub.
Record your change first (the \cmd{git add} and \cmd{git commit} commands above), then run
\cmd{git pull} again.

\subsection{End a session}

\step{Close the window}
When the prompt is showing, type \code|exit| and press \key{Enter}, or close the window with
the \ui{X} button. Nothing needs to be saved or switched off; the environment is simply
activated again by the opening sequence next time.

\begin{caution}[Do not update the software before the hand-in]
The versions installed by this guide are known to work together. Until the project has been
handed in on December~1, decline offers to upgrade Python packages (including the notice that a
new version of pip is available), and do not run commands that upgrade them. A working
installation is worth more than a current one.
\end{caution}
